\section{Integration schemes}

In order to integrate the homogeneous differential equations of motion
we discretize time as $t_{i}=t_{0}+i\Delta t$ with $i=1,\ldots,n_{steps}$,
where $n_{steps}$ denotes the number of sub-intervals. In JETSPIN
three different integration schemes: the first-order accurate Euler
scheme, the second-order accurate Heun scheme (sometimes called second-order
accurate Runge-Kutta), and the fourth-order accurate Runge-Kutta scheme\cite{press2007numerical}.
The user can select a specific scheme by using appropriate keys in
the input file, as described in Sec \ref{sec:Input-file}. If the
airdrag is activated, the explicit strong order scheme by Platen \cite{platen1987derivative,kloeden1992numerical,platen2010numerical}
have to be selected, whereof the order of strong convergence was evaluated
equal to 1.5.

This scheme avoids the use of derivatives by corresponding finite
differences in the same way as Runge-Kutta schemes do for ODEs in
a deterministic setting, and it is briefly summarized as follows.

Let us consider a Brownian motion vector process $\textbf{X}=\left\{ \textbf{X}_{t},t\right\} $
of \textit{d}-dimensional satisfying the stochastic differential equation

\begin{equation}
\frac{d\textbf{X}}{dt}=\textbf{a}\left(t,X^{1},\ldots,X^{d}\right)+\textbf{b}d\Omega\label{eq:generic-sde}
\end{equation}
where $\textbf{a}$ and $\textbf{b}$ are vectors of \textit{d}-dimensional
usually called drift and diffusion vector coefficients, and $\Omega\left(t\right)$
denotes a Wiener process. Denoted $Y_{t}^{k}$ the approximation for
the \textit{k}-th component of the vector $\textbf{X}$ at time $t$,
the integrator has the following form:

\begin{equation}
Y_{t+\varDelta t}^{k}=Y_{t}^{k}+b^{k}\Delta\Omega+\frac{1}{2\sqrt{\Delta t}}\left[a^{k}\left(\widetilde{\boldsymbol{\Upsilon}}_{+}\right)-a^{k}\left(\widetilde{\boldsymbol{\Upsilon}}_{-}\right)\right]\Delta\varPsi+\frac{1}{4}\left[a^{k}\left(\widetilde{\boldsymbol{\Upsilon}}_{+}\right)+2a^{k}+a^{k}\left(\widetilde{\boldsymbol{\Upsilon}}_{-}\right)\right]\Delta t,\label{eq:integrator-platen}
\end{equation}

with the vector supporting values

\begin{equation}
\begin{alignedat}{1}\widetilde{\boldsymbol{\Upsilon}}_{\pm} & =\textbf{Y}_{t}+\textbf{a}\Delta t\pm\textbf{b}\sqrt{\Delta t}\\
\widetilde{\boldsymbol{\Phi}}_{\pm} & =\widetilde{\boldsymbol{\Upsilon}}_{+}\pm\textbf{b}\left(\widetilde{\boldsymbol{\Upsilon}}_{+}\right)\sqrt{\Delta t}.
\end{alignedat}
\end{equation}

Here, $\Delta\Omega$ and $\Delta\varPsi$ indicate normally distributed
random variables constructed from two independent $N\left(0,1\right)$
standard Gaussian distributed random variables ($U_{1}$, $U_{2}$)
by means of the following linear transformation:

\begin{equation}
\begin{alignedat}{1}\Delta\Omega & =U_{1}\sqrt{\Delta t}\\
\Delta\varPsi & =\frac{1}{2}\Delta t^{3/2}\left(U_{1}+\frac{1}{\sqrt{3}}U_{2}\right).
\end{alignedat}
\label{eq:random-var-1}
\end{equation}

Note that in JETSPIN the time step $\Delta t$ is automatically rescaled
by the quantity $\tau$, in accordance with the mentioned dimensionless
scaling convention.
